            Parameter {\bf out} is the cmemory or rmemory mode output variable.

            
            In the case of Zynq the interface uses AXI bus GP0.            
            The size of the output data type must not exceed 256 bits as the CPU
            interface software limits burst transfers to 8 32-bit words.
            
            Boards having a serial interface to the CPU use that interface.
            The size of the output data type is not limited.
            
            This interface can use any memory port and that port can also be
            used by other application code, but no arbitration of shared access
            via a port is provided. Coincident access will usually result in
            data corruption. If it is not feasible to dedicate a port to CPU
            access then the CPU application code must limit acccess to the
            memory to times when it is known that the FPGA logic will not also
            be accessing the memory port.
            
            On running info2c on the .json file the following C function will have been
            defined for interface {\em name} -
            
            \begin{tabular}{| l | l |}
            \hline
            function                   &                              \\ \hline \hline
            fpga\_{\em name}\_write()  & write a datum to FPGA memory \\ \hline
            \end{tabular}
            
            Example: \\
            If 3PL source file prog.3pl contains the following \\
            \begin{lstlisting}[gobble=12, language=threepl]
               class(cpu, cpu_api())
               rmemory(m, 2, "uint:8", "uint:16");
               .
               .
               cpu.write_memory(m <- "MEM2", 1);
               .
               .
            \end{lstlisting}
            then for zynq the CPU code should contain something like \\
            \begin{lstlisting}[gobble=12, language=C]
                   #include "prog.h"
                   .
                   .
                   fpga_t *fpga = NULL;
                   int     d;
                   .
                   .
                   fpga_open(portName, INTERRUPTS);
                   .
                   .
                   fpga_MEM2_write(fpga, addr, d);
                   .
            \end{lstlisting}
            where addr is the memory addressd the data to be written.

            This procedure creates no in-line target code so the current clock domain
            when this procedure is called is not relevant. The clock domain is not
            changed by this procedure,
